%% ============================================================
%%  REFERINȚE
%% ============================================================
\clearpage
\thispagestyle{empty}
\addcontentsline{toc}{chapter}{Referințe}
\vspace*{1cm}
{\noindent\color{onm}\rule{\linewidth}{3pt}}\par\vspace{3mm}
{\noindent\fontsize{28}{32}\selectfont\bfseries\color{ink} Referințe}\par\vspace{3mm}
{\noindent\color{onm}\rule{\linewidth}{3pt}}\par\vspace{8mm}

\noindent Volumul de față este autoconținut; demonstrațiile, definițiile și exemplele nu trimit la resurse externe. Lista de mai jos cuprinde lucrările la care autorul s-a raportat în procesul de redactare --- fie pentru a verifica formulări, fie ca surse de inspirație pentru structura prezentării.

\vspace{6mm}
\begin{thebibliography}{9}

\bibitem{exc11vol1}
Marius Perianu, Ioan Balică, Dan Zaharia,
\textit{Matematică de excelență pentru clasa a XI-a, Vol.\ I: Algebră},
Editura Paralela 45, Pitești, 2012.

\bibitem{exc11vol2}
Marius Perianu, Ioan Balică, Dan Zaharia,
\textit{Matematică de excelență pentru clasa a XI-a, Vol.\ II: Analiză},
Editura Paralela 45, Pitești, 2012.

\bibitem{exc12vol1}
Marius Perianu, Ioan Balică, Dan Zaharia,
\textit{Matematică de excelență pentru clasa a XII-a, Vol.\ I: Algebră},
Editura Paralela 45, Pitești, 2013.

\bibitem{exc12vol2}
Marius Perianu, Ioan Balică, Dan Zaharia,
\textit{Matematică de excelență pentru clasa a XII-a, Vol.\ II: Analiză},
Editura Paralela 45, Pitești, 2013.

\bibitem{titu}
Titu Andreescu, Dorin Andrica,
\textit{Algebră liniară, culegere de probleme},
GIL, Zalău, 2004.

\end{thebibliography}

